% Non-CM and integral-model part. Uses only standard LaTeX, amsmath/amssymb,
% and the common macros authorized in the main paper.
\section{Non-CM sources and the integral-model comparison}

\subsection{THETA3-C01: the precise congruence $H^1$ obstruction}\label{C01}

\paragraph{Statement.}
Let $X$ be a smooth projective congruence unitary ball quotient in the
setup of Y.~Liu, \emph{Fourier--Jacobi cycles and arithmetic relative
trace formula}, with Hermitian rank $m\geq3$. If $B$ is a simple complex
abelian variety which is not of CM type, then
\[
 \Hom_{\mathrm{Hdg}}\bigl(H^1(B,\Q),H^1(X,\Q)\bigr)=0.
\]
The same conclusion holds for any character-generated $H^1$ subspace
which is a sum of images of CM Hodge structures. Sums, tensor products,
duals, and Hodge subquotients of such structures remain of CM type.

\paragraph{Proof.}
Liu's Corollary~4.20, pp.~54--55, gives the CM isogeny decomposition of
the whole Albanese in rank $m\geq3$. Consequently $H^1(X,\Q)$ has
commutative Mumford--Tate group. A subquotient of a CM Hodge structure
is CM, because its Hodge homomorphism factors through an image of the
same torus. A nonzero Hodge morphism from $H^1(B,\Q)$ is injective:
its kernel is a rational sub-Hodge structure, and simplicity follows
from the equivalence between abelian varieties up to isogeny and
polarizable rational Hodge structures of weight one. Such an
injection into a CM structure would force $B$ to be CM, a
contradiction. A finite collection of CM structures is controlled by
a product of tori; all its tensor constructions and subquotients are
therefore controlled by a torus as well. $\square$

This is a deduction from the cited CM decomposition, not a new proof
of that decomposition. Liu's Hermitian rank is $m$, so the ball has
dimension $m-1$. In rank two the cited statement concerns the
endoscopic Albanese part, not the entire Albanese. The attached
van Hoften--Sempliner paper gives the relevant rank-three result in
Proposition~3.1.7, Theorem~3.1.11, and \S3.1.12. Increasing the ball
dimension does not remove this $H^1$ obstruction. Neither statement
applies to arbitrary noncongruence subgroups or to all higher
cohomology of the ball quotient.

\paragraph{Framework anchor.}
The relevant realization is the Albanese, not an arbitrary higher-degree
Hodge summand.  Grothendieck's letter of 12 August 1964, cote~44,
scans~7--8, explicitly places $\operatorname{Pic}^0$ and
$\operatorname{Alb}^0$ in the correspondence and polarization framework.
The CM restriction above is the additional automorphic input, not a
statement extracted from that letter.

\subsection{THETA3-C02: a trace-zero selection rule on covers}\label{C02}

\paragraph{Statement.}
Let $f:X\to S$ be a finite \'etale cover of connected smooth
projective complex surfaces of degree $d$. Then
\[
 H^1(X,\Q)=f^*H^1(S,\Q)\oplus
       \ker\bigl(f_*:H^1(X,\Q)\to H^1(S,\Q)\bigr).
\]
If $u$ belongs to the second summand and $v_1,v_2,v_3\in H^1(S,\Q)$,
then
\[
 \int_Xu\smile f^*v_1\smile f^*v_2\smile f^*v_3=0.
\]
If $H^1(S,\Q)$ is CM, every simple non-CM sub-Hodge structure of
$H^1(X,\Q)$ belongs to this trace-zero summand.

\paragraph{Proof.}
The transfer satisfies $f_*f^*=d\,\mathrm{id}$. Therefore
$d^{-1}f^*f_*$ is a projector with image $f^*H^1(S,\Q)$ and kernel
$\ker f_*$. The projection formula gives
\[
 \int_Xu\smile f^*v_1\smile f^*v_2\smile f^*v_3
 =\int_S f_*u\smile v_1\smile v_2\smile v_3=0.
\]
The restriction of $f_*$ to a simple non-CM substructure vanishes by
the argument of C01. $\square$

Thus one new non-CM factor and three old theta factors cannot give
the desired nonzero four-form period. Two new factors, for example a
holomorphic form and its conjugate, are the first possible case.

\paragraph{Framework anchor.}
The maps here are the pullback and trace of an actual finite geometric
correspondence.  Their action on the Albanese realization fits cote~44,
scans~7--8.  The vanishing calculation is a projection-formula statement;
it is not an assertion that arbitrary projectors in a realization are
represented by algebraic cycles.

\subsection{THETA3-C03: a mixed non-CM and CM period}\label{C03}

\paragraph{Hypotheses and source.}
Let $S_0=\Gamma_0\backslash D$ be a compact principal congruence
ball surface in the exact setup of 032, \S4. Suppose a compact ball
surface $X_0$ commensurable with $S_0$ has a nonconstant morphism
$a_0:X_0\to E_0$ to a non-CM elliptic curve.
Schoen's \emph{An arithmetic ball quotient surface whose Albanese
variety is not of CM type}, pp.~133--134, supplies such a
commensurability class over $\Q(\zeta_5)/\Q(\sqrt5)$, with
$j(E_0)=2^{11}$. To use it, choose 032's Hermitian space to be the
particular rank-three space in Schoen's construction. Multiplication
of the form by an element of the totally real field with suitable
real signs gives signatures $(2,1)$ and $(3,0)$. Equality of CM
fields alone would not imply commensurability. The determinant image
of an arithmetic unitary subgroup is finite: it is commensurable
with an integral subgroup, and integral norm-one units in a CM field
form a finite group. A finite-index subgroup therefore lies in
$\SU(V)$. Keep the principal congruence surface as the source of the
theta construction.

\paragraph{Statement.}
Fix an allowed CM sign function $v$ of 032 with $v(1)=1$.
There exist a sufficiently deep principal cover $S$ of $S_0$, a
common finite \'etale cover $X$ of $X_0$ and $S$, and a theta
one-form $\beta$ on $S$ in the complex span of rational source maps
from $U(v)$ such that, for a nonzero invariant form $\alpha_0$ on
$E_0$,
\[
 \xi=a^*\alpha_0\wedge f^*\beta\ne0,
 \qquad \int_X\xi\wedge\overline\xi>0,
\]
where $a:X\to E_0$ and $f:X\to S$ are the induced maps. Moreover
there are a CM abelian variety $B(v)$, a morphism $b:X\to B(v)$,
and a form $\gamma\in H^0(B(v),\Omega^1)$ with
$a^*\alpha_0\wedge b^*\gamma\ne0$. The image of $(a,b)$ has
dimension two.

\paragraph{Proof.}
The pullback $\alpha$ of $a_0^*\alpha_0$ to $D$ is nonzero. Choose
$x\in D$ with $\alpha_x\ne0$. By 032 Proposition~4.5,
pp.~21--22, finitely many actual finite Schwartz inputs give theta
one-form values spanning $T_x^*D$. At least one resulting $\beta_x$
is independent of $\alpha_x$. These finitely many inputs descend to
a sufficiently deep principal cover $S$ of $S_0$. Refine it using
032 Theorem~5.1, p.~26, and Corollary~6.7, p.~45, so that $\beta$
has its rational CM-source expansion. Only now form $X$ from the
intersection of the torsion-free commensurable lattices defining
$X_0$ and $S$. The common finite cover preserves the nonzero wedge
at a lift of $x$.

On the projective surface $X$, a nonzero holomorphic two-form has
strictly positive integral of $\xi\wedge\overline\xi$. Indeed, in
local coordinates $\xi=g\,dz_1\wedge dz_2$, this is a nonnegative
real volume form and is positive wherever $g\ne0$.

The CM extraction writes $[\beta]$ as a finite complex sum of
$h_j(u_j)$, where $h_j:U(v)\to H^1(S,\Q)$ are rational Hodge maps
and $u_j$ lie in the selected $(1,0)$ line. Holomorphic one-forms on
a compact K\"ahler manifold are the unique holomorphic
representatives of their classes, so the same expansion holds for
the forms. At least one term satisfies
$\alpha\wedge f^*h_j(u_j)\ne0$. By the Albanese realization used
in 032 Lemma~2.2, a positive integer multiple of $h_j$ is induced
by an algebraic map $S\to B(v)$. Composing with $f$ gives $b$.
If $(a,b)$ had image dimension at most one, its differential would
have rank at most one on a dense open set, forcing this wedge to
vanish. Therefore its image has dimension two, and
\[
 (a,b)_*[X]\in
 \CH^{\dim B(v)-1}(E_0\times B(v))_{\Q}
\]
is an algebraic image-surface cycle with the stated nonzero pairing.
$\square$

This construction uses two new factors, $\alpha$ and
$\overline\alpha$, and two CM theta factors. It does not realize
the new forms as adelic character theta lifts on a noncongruence
subgroup. It also does not identify a previously unknown Hodge class
as algebraic; C05 explains the limited Hodge content of this target.

\paragraph{Framework anchor.}
The surface image is an actual algebraic correspondence, and the nonzero
fourfold period is its Betti pairing.  Cote~44, scans~7--8, is the
correspondence/polarization anchor for that passage.  The extra non-CM
Albanese factor and the transverse theta form are supplied by the two
modern source constructions, not by an archival universality statement.

\subsection{THETA3-C04: the exact projector and rank gate}\label{C04}

\paragraph{Statement.}
Let $A$ be smooth projective, and let
$T\subset H^{2p}(A,\Q)(p)$ be a $d$-dimensional space of rational
Hodge classes. Suppose an algebraic correspondence acts as a
projector $e$ onto $T$. For algebraic cycles $Z_1,\ldots,Z_r$ of
codimension $p$, a basis $t_1,\ldots,t_d$ of $T$, and a
nondegenerate rational bilinear form $q_T$ on $T$, put
\[
 M_{ij}=q_T\bigl(e\,\operatorname{cl}(Z_i),t_j\bigr).
\]
If $\operatorname{rank}M=d$, every element of $T$ is algebraic.
Alternatively, if an algebra $\mathcal D$ of algebraic
correspondences preserves $T$, it suffices that the
$\mathcal D$-orbit of $e\operatorname{cl}(Z)$ spans $T$ for one
cycle $Z$. If $T$ is a simple $\mathcal D$-module, a nonzero
$e\operatorname{cl}(Z)$ suffices.

\paragraph{Proof.}
Every $e\operatorname{cl}(Z_i)$ is algebraic. Nondegeneracy
identifies $T$ with its dual, so full row rank means that these
vectors span $T$. In the orbit version, the finite-dimensional
orbit span is generated by algebraic correspondences applied to an
algebraic cycle. A nonzero submodule of a simple module is the whole
module. $\square$

A single nonzero scalar pairing does not give this rank assertion.
For example, $\Q(1,0)\subset\Q^2$ pairs nontrivially with $(1,1)$
under the ordinary dot product, although $(1,1)$ is outside that
subspace. This is a linear-algebra counterexample to an inference,
not an example of a nonalgebraic Hodge class. The pure embedding-line
divisors of 032 Proposition~3.1 supply the necessary CM extraction
in that paper. For a general non-CM $H^1$, the polarization provides
a sum of rank-one tensors; it does not supply individual
holomorphic eigenline projectors as algebraic correspondences.

\paragraph{Framework anchor.}
The distinction between an inverse or projector on cohomology and one
induced by an algebraic correspondence is explicit in cote~16,
scan~44, conditions~(iv) and~(iv bis), with the discussion continuing on
scans~45--46.  This finding isolates the same existence requirement and
adds the finite-dimensional rank test needed for the specified target.

\subsection{THETA3-C05: one non-CM elliptic factor gives a known consequence}\label{C05}

\paragraph{Statement; no novelty claim.}
Let $E_0$ be a non-CM elliptic curve, and let $B$ be a CM abelian
variety for which the rational Hodge conjecture holds. Then every
Hodge class on $E_0\times B$ is algebraic. More precisely,
\[
 \Hdg^p(E_0\times B)
 =1\otimes\Hdg^p(B)
 \oplus[\mathrm{pt}]\otimes\Hdg^{p-1}(B),
\]
where $\Hdg^p(Y)$ denotes the rational $(0,0)$ classes in
$H^{2p}(Y,\Q)(p)$.

\paragraph{Proof.}
The classical Hodge-group facts are
$\operatorname{Hg}(E_0)=\operatorname{SL}(H^1(E_0,\Q))$ and
$\operatorname{Hg}(B)=T$ for a torus $T$. For the first statement,
see Moonen, \emph{Notes on Mumford--Tate groups}, \S2,
Example~2.6, p.~11. The connected Hodge group $H$ of the product
projects onto both factors. Its derived group projects onto
$\operatorname{SL}_2$ and trivially to $T$. As a subgroup of
$\operatorname{SL}_2\times1$ with full projection to
$\operatorname{SL}_2$, it equals that group. The remaining quotient
surjects onto $T$, whence $H=\operatorname{SL}_2\times T$.
The only $\operatorname{SL}_2$ invariants in
$H^*(E_0,\Q)=\Q\oplus H^1(E_0,\Q)\oplus\Q(-1)$ occur in degrees
zero and two. Taking invariants in the K\"unneth decomposition gives
the formula. The two kinds of class are represented by a pullback of
a cycle on $B$, or by a point on $E_0$ times such a cycle. $\square$

\paragraph{Framework anchor.}
This is a consequence inside the tensor category generated by the
weight-one realizations, with the elliptic degree-two factor represented
by a divisor.  It uses the Albanese/correspondence framework of cote~44,
scans~7--8; it does not supply the additional higher-degree algebraic
correspondences discussed in cote~16, scans~43--46.

\subsection{The precise unramified local normalization}

Let $k_0/\Q_p$ be finite with $p$ odd, residue cardinality $q$, and
uniformizer $\pi$. Let $k/k_0$ be unramified quadratic and
$\mathcal O=\mathcal O_k$. We use Li--Zhang's principally polarized
unitary Rapoport--Zink space $\mathcal N_n$ of signature $(1,n-1)$,
and its nonsplit rank-$n$ Hermitian space $\mathbb V_n$ of special
quasi-homomorphisms. For a full-rank lattice
$L=\mathcal O x_1+\cdots+\mathcal O x_n\subset\mathbb V_n$, define
\[
 \operatorname{Int}(L)=\chi\left(\mathcal N_n,
  \mathcal O_{\mathcal Z(x_1)}\otimes^{\mathbf L}\cdots
                    \otimes^{\mathbf L}\mathcal O_{\mathcal Z(x_n)}\right).
\]
This is a derived Euler characteristic, not a definition by the
ordinary scheme-theoretic intersection length. The source proves
basis independence. Li--Zhang, \emph{Kudla--Rapoport cycles and
derivatives of local densities}, Theorem~1.2.1, equivalently
Theorem~3.4.1, identifies it with
\[
 \operatorname{Int}(L)=-\left.\frac{d}{dX}
                           \operatorname{Den}(X,L)\right|_{X=1}.
\]
Their Theorem~3.5.1 gives the exact polynomial
\begin{align*}
 \operatorname{Den}(X,L)
 &=\sum_{L\subset N\subset N^\vee}
      X^{2\operatorname{length}_{\mathcal O}(N/L)}m(t(N);X),\\
 m(t;X)&=\prod_{i=0}^{t-1}\bigl(1-(-q)^iX\bigr).
\end{align*}
Here $t(N)$ is the number of positive elementary divisors. The
nonsplit space in this convention has odd determinant valuation.
The interpolation normalization is
\[
 \operatorname{Den}((-q)^{-k},L)
 =\frac{\operatorname{Den}(\langle1\rangle^{n+k},L)}
        {\operatorname{Den}(\langle1\rangle^{n+k},\langle1\rangle^n)}.
\]
The identity $\operatorname{Int}=-\operatorname{Den}'(1)$ has no
omitted logarithm. Conversion to a Whittaker derivative introduces
$\log(q^2)$ and explicit unramified factors as in the source's
Remark~3.4.2. The following calculations use this precise
normalization. Author-PDF locators are \S1.2, pp.~3--4;
Theorem~3.4.1, p.~17; Theorem~3.5.1, p.~18; and
Corollary~3.5.3, p.~19.

\subsection{THETA3-C06: an explicit negative derived intersection}\label{C06}

\paragraph{Statement.}
For $n=3$, let $L$ have Hermitian Gram matrix $\pi I_3$.
Then in the specified nonsplit local model,
\[
 \operatorname{Den}(X,L)=(1-X)\bigl(1+(q-q^2)X+X^2\bigr),
 \qquad \operatorname{Int}(L)=q+2-q^2<0.
\]
In particular the value at $q=3$ is $-4$.

\paragraph{Proof.}
The determinant valuation is three, hence odd, so the ambient
Hermitian space is the required nonsplit $\mathbb V_3$. Every
integral overlattice $N\supset L$ has
\[
 \operatorname{val}(N)=3-2\operatorname{length}_{\mathcal O}(N/L)
 \geq0.
\]
Thus its index length is zero or one. The length-zero term is $L$,
of type three. Since $L^\vee=\pi^{-1}L$, a length-one overlattice
has the form $N=L+\mathcal O\pi^{-1}v$, where $v\bmod\pi$ spans
a line in $L/\pi L$. Its pairings with $L$ are integral. Its
remaining integrality condition is
$h(v,v)\in\pi^2\mathcal O$, equivalently that the residual line
is isotropic for the nondegenerate Hermitian form $I_3$ over
$\mathbb F_{q^2}/\mathbb F_q$. Distinct residual lines give
distinct overlattices. Such an $N$ has determinant valuation one,
so its elementary divisors are $(0,0,1)$ and $t(N)=1$.

Here is a direct count of these lines. Let $a_r$ count vectors in
$\mathbb F_{q^2}^r$ whose coordinate norms sum to zero. The norm
map has one preimage of zero and $q+1$ preimages of each nonzero
element, so
\[
 a_1=1,\qquad
 a_r=(q+1)q^{2r-2}-q a_{r-1}.
\]
Consequently $a_2=q^3+q^2-q$ and $a_3=q^5-q^3+q^2$. Dividing the
nonzero isotropic-vector count by $q^2-1$ gives
\[
 \frac{q^5-q^3+q^2-1}{q^2-1}=q^3+1
\]
isotropic lines. The density formula therefore becomes
\begin{align*}
 \operatorname{Den}(X,L)
 &=(1-X)(1+qX)(1-q^2X)+(q^3+1)X^2(1-X)\\
 &=(1-X)\bigl(1+(q-q^2)X+X^2\bigr).
\end{align*}
Differentiating and invoking Li--Zhang gives
$\operatorname{Int}(L)=q+2-q^2$. Since
$q^2-q-2=(q-2)(q+1)>0$ for odd $q\geq3$, it is negative.
$\square$

This is a specialization of the source theorem, not a claim of a
new local-density theorem. It is consistent with the independent
horizontal and vertical calculation in Li--Zhang's proof of
Lemma~6.2.1, pp.~35--36: at $a_3=1$ the two terms are
$(q+1)+(1-q^2)$. A separate finite-field script in this run counted
$28,126,344,1332$ lines for $q=3,5,7,11$, respectively, and obtained
intersection values $-4,-18,-40,-108$. These finite checks do not
replace the general proof.

Thus a positive definite Hermitian norm cannot be identified with
an arbitrary local arithmetic intersection without controlling the
derived signs and the chosen cycle combination. Positivity of
representation densities at their interpolation arguments also
does not imply positivity of the derivative at one.

\paragraph{Framework anchor.}
The precise intersection-theoretic analogue is SGA~4$\frac12$,
\emph{Cycle}, Remark~2.3.9, p.~150, where the multiplicity is an
alternating Tor sum.  The negative example prevents identifying this
derived local number with the positive polarized pairing of cote~44,
scans~7--8, without an additional comparison theorem.

\subsection{THETA3-C07: a positive local sector with an exact bound}\label{C07}

\paragraph{Statement.}
Let $M$ be a self-dual Hermitian $\mathcal O$-lattice of rank $n-1$,
and let $x\perp M$ satisfy $h(x,x)=\pi^{2a+1}$ with $a\geq0$.
For $L=M\perp\mathcal O x$,
\[
 \operatorname{Den}(X,L)=(1-X)\sum_{j=0}^aX^{2j},
 \qquad \operatorname{Int}(L)=a+1\geq1.
\]

\paragraph{Proof.}
An integral overlattice $N\supset L$ splits as
$M\perp(N\cap M_k^\perp)$. Indeed, for $y\in N$, its orthogonal
projection to $M_k$ belongs to $M$ because $h(y,M)\subset\mathcal O$
and the pairing on $M$ is perfect. Subtract this projection.
The complementary lattice is $\pi^{-j}\mathcal O x$. Containment
and integrality give $j\geq0$ and $2a+1-2j\geq0$, hence
$0\leq j\leq a$. Each overlattice has index length $j$ and type
one, since its remaining positive elementary divisor is the odd
integer $2a+1-2j$. The density formula is the displayed sum.
Differentiation at one and the local theorem give $a+1$.
$\square$

The rank-one deformation interpretation after cancellation of the
self-dual part appears in Li--Zhang \S2.11 and Examples~2.4.3 and
3.5.2. The preceding proof does not require that interpretation.
It gives a normalized positive local certificate; comparison with
a different archimedean mixed period still requires a matching
theorem for its test data and factors.

\paragraph{Framework anchor.}
Here the exact local model permits a positive calculation, while the
cycle-class and intersection comparison still has the hypotheses of
SGA~4$\frac12$, \emph{Cycle}, Theorem~2.3.8 and Remark~2.3.9,
pp.~148--150.  Its connection to the correspondence/polarization problem
of cote~44, scans~7--8, requires the geometric identification in C17.

\subsection{THETA3-C08: vertical positivity does not detect a generic class}\label{C08}

\paragraph{Statement and counterexample.}
A positive ordinary local intersection on an integral model need
not detect a nonzero generic-fibre cycle. Let $R$ be a discrete
valuation ring with uniformizer $\pi$, and consider $\mathbb P^1_R$.
Its special fibre $F$ and a section $H$ meet with local intersection
ring $R/(\pi)$ and length one. Nevertheless $F$ has empty generic
fibre, so its restriction to the generic Chow group and cohomology
is zero. $\square$

An actual horizontal flat cycle of the required relative
codimension has an algebraic generic fibre by construction. The
remaining issue is identifying its class with the desired Hodge
tensor. The specialization criterion below makes the available
functoriality precise; it does not construct the missing cycle.

The basic-locus uniformization retains a geometric factor as well.
Vollaard--Wedhorn, \emph{The supersingular locus of the Shimura
variety of $\operatorname{GU}(1,n-1)$ II}, equation~(5.4.2), p.~34,
is
\[
 I(\Q)\backslash
 \bigl(\mathcal N_{\mathrm{red}}\times G(\mathbb A_f^p)/K^p\bigr)
 \simeq \mathcal M^{\mathrm{ss}}.
\]
It does not identify the whole supersingular locus with a finite
Shimura set. In rank three the Rapoport--Zink factor contributes
Fermat curves, as recalled in the source's introduction, p.~4.
Definite mass formulas can describe finite indexing data, while
local intersections depend on the retained geometry.

Li--Zhang's Theorem~4.2.1, p.~21, decomposes the horizontal part of
a rank-$(n-1)$ special cycle into primitive quasi-canonical lifting
cycles indexed by integral overlattices of type at most one.
Theorem~6.1.3, p.~34, compares the horizontal intersection function
with the corresponding horizontal density derivative. These are
more refined existing inputs than an unqualified positivity claim
for the total intersection.

\paragraph{Framework anchor.}
Cote~16, scans~44--46, explicitly discusses specialization in the
standard-conjecture framework.  The example above explains why a
vertical local intersection is not, by itself, a generic realization
class.  C17 states the actual relative-correspondence hypotheses under
which the cited SGA specialization comparison applies.

\subsection{THETA3-C09: finite lifting valuations and the infinite-product obstruction}\label{C09}

\paragraph{Statement.}
For $e_s=q^{s-1}(q+1)$,
\[
 \sum_{s\geq1}\frac1{e_s}=\frac{q}{q^2-1}.
\]
If $t_s\in\C_p$ have normalized valuations $v(t_s)=1/e_s$, the
infinite product $\prod_{s\geq1}t_s$ does not converge in $\C_p$.

\paragraph{Proof.}
The sum is geometric. The valuations of its partial products tend
to the displayed finite value, so the absolute values tend to a
nonzero number. A convergent sequence of these products would have
a nonzero limit, forcing its successive ratios $t_s$ to tend to
one. But $v(t_s)>0$ implies $|t_s|<1$ and $|1-t_s|=1$ for every
$s$. This is impossible. $\square$

The relevant primary source is Gross, \emph{On canonical and
quasi-canonical liftings}, Proposition~5.3(2),(3), manuscript p.~9,
proof pp.~9--11, in the original MPI/SFB~85-30 scan appended to
Gross--Zagier. These are PDF pages~77--78 of the retrieved combined
scan. In the unramified quadratic case it gives the degree
$[W_s:W]=e_s$ and valuation $v_{W_s}(t_s)=1$ of the formal modulus.
The original scan was visually checked. Li--Zhang \S4.1, p.~20,
uses the same extension degree $q^s(1+q^{-1})$.

Normalize the base valuation by $v(\pi)=1$. Then
$v(t_s)=1/e_s$. The canonical modulus $t_0\in\pi W$ satisfies
$v_{W_s}(t_s-t_0)=1$. In $W[[t]]$ the intersection with the canonical
section is therefore
\[
 W_s/(t_s-t_0),
\]
of ordinary length one over $W_s$ and over $W$, since the residue
fields agree. The branch degree $e_s$, its local intersection
length one, and the base-normalized coordinate valuation $1/e_s$
are different invariants. For finite $S$,
\[
 v\left(\prod_{s=1}^St_s\right)
 =\frac{q(1-q^{-S})}{q^2-1},\qquad
 v\left(\prod_{s=1}^S
             \operatorname{Norm}_{W_s/W}(t_s)\right)=S.
\]
The norm uses complete conjugate orbits and has valuation one at
each level. None of these finite identities is itself an identity
for a theta section. A comparison by divisors must specify the
section, its integral frame, and which branches or norms occur.
For the comparison with C11 over $\Q_p$, set this $q=p$; the
Lubin--Tate residue cardinality in C11 is instead $Q=p^2$.

\paragraph{Framework anchor.}
Gross's deformation divisors supply a concrete local intersection problem.
The exact normalization matters before applying the intersection and
cycle-class formalism of SGA~4$\frac12$, \emph{Cycle}, pp.~148--150.
The correspondence viewpoint in cote~44, scans~7--8, does not identify
these local divisors with a chosen analytic theta section.

\subsection{THETA3-C10: non-CM arithmetic theta and its realization}\label{C10}

\paragraph{Elementary observation.}
If an algebraic cycle $Z$ is homologically trivial, every ordinary
cohomological pairing with $Z$ is zero. A nonzero arithmetic height
may detect information about rational equivalence but does not
supply a nonzero ordinary Hodge class.

\paragraph{Proof.}
Homological triviality is precisely
$\operatorname{cl}_{\mathrm B}(Z)=0$. Its cup product with every
Betti class is zero, and the same statement holds in the compatible
\'etale realization. A height does not change this vanishing.
$\square$

There is an existing non-character extension relevant to this
boundary. Li--Liu, \emph{Chow groups and $L$-derivatives of
automorphic motives for unitary groups, II}, Example~1.10,
pp.~6--7, includes the representation attached to
$\operatorname{Sym}^{2r-1}$, $r\geq2$, of a non-CM elliptic curve,
under its modularity and local hypotheses. Good reduction at all
nonsplit finite places gives $S_\pi=\varnothing$; when
$r[F:\Q]$ is odd the stated case applies. Theorem~1.4 produces a
nonzero localized $\CH^r_{\mathrm{hom}}$ from a simple central zero,
with the technical Hypothesis~6.6 removable in this empty-$S_\pi$
case by Remark~1.9. Theorem~1.5's height formula also requires the
stated generating-series modularity hypothesis. These cited source
results are carried as background, not reproved here.

Thus non-CM input is available in arithmetic theta theory, but this
output belongs to homologically invisible cycles and
Abel--Jacobi or motivic extension data. A non-CM Hodge-class problem
still requires an algebraic correspondence with the relevant
nonzero realization and the rank condition C04.

\paragraph{Framework anchor.}
The output in this source is homologically trivial, so its nonvanishing
height does not give a nonzero class in the pure cohomological realization.
This is the precise distinction from the algebraic cohomology classes
and correspondence inverses discussed in cote~16, scans~43--46.

\subsection{THETA3-C11: the exact local derivative slope, with sharpness}\label{C11}

\paragraph{Statement.}
Let $p$ be odd, put $Q=p^2$, and choose $\pi\in\C_p$ with
$v_p(\pi)=1$. Set
\[
 \ell(T)=\sum_{r\geq0}\frac{T^{Q^r}}{\pi^r},\qquad
 E(Z)=\ell^{-1}(Z),\qquad D=\ell'(T)^{-1}\frac{d}{dT}.
\]
For every $f\in\mathcal O_{\C_p}[[T]]$ and every integer $a\geq1$,
\begin{align*}
 v_p(D^af(0))
 &\geq v_p(a!)-\frac{a-1}{p^2-1}\\
 &=\frac{ap}{p^2-1}+\frac1{p^2-1}
                     -\frac{s_p(a)}{p-1},
\end{align*}
where $s_p(a)$ is the base-$p$ digit sum and $v_p(0)=+\infty$.
Moreover $D^p=0$ after reduction modulo $p$.
The bound is attained infinitely often: for $f(T)=T$ and $a=Q^r$,
\[
 v_p(D^{Q^r}T(0))=\frac{(Q^r-1)p}{p^2-1}.
\]

\paragraph{Proof of the bound.}
Choose $\beta$ with $\beta^{Q-1}=\pi$, so
$v_p(\beta)=1/(Q-1)$. The scaled logarithm is
\[
 L(T)=\frac{\ell(\beta T)}\beta
 =\sum_{r\geq0}
     \pi^{(Q^r-1)/(Q-1)-r}T^{Q^r}.
\]
It is integral and has linear coefficient one, since for $r\geq1$
\[
 \frac{Q^r-1}{Q-1}-r=1+Q+\cdots+Q^{r-1}-r\geq0.
\]
Recursive formal inversion over any commutative ring therefore gives
\[
 U(T)=\frac{E(\beta T)}\beta
       \in T\mathcal O_{\C_p}[[T]],\qquad U'(0)=1.
\]
Integrality of $f$ implies
\[
 f(E(\beta T))-f(0)=f(\beta U(T))-f(0)
                         \in\beta T\mathcal O_{\C_p}[[T]].
\]
The chain rule in the logarithmic coordinate gives
\[
 \frac{d}{dZ}f(E(Z))=(Df)(E(Z)),\qquad
 f(E(Z))=\sum_{a\geq0}\frac{D^af(0)}{a!}Z^a.
\]
Its coefficient after substituting $Z=\beta T$ is
$\beta^aD^af(0)/a!$, of valuation at least $v_p(\beta)$ for
$a\geq1$. This proves the bound. Legendre's formula
$v_p(a!)=(a-s_p(a))/(p-1)$ gives the second expression. Since
$s_p(a)\leq(p-1)(1+\lfloor\log_pa\rfloor)$, the error from the
linear term is explicitly logarithmic.

\paragraph{Proof of nilpotence.}
We have
\[
 \ell'(T)=\sum_{r\geq0}(Q/\pi)^rT^{Q^r-1}\equiv1\pmod p.
\]
Thus $D$ reduces to $d/dT$ over $\mathcal O_{\C_p}/p$.
The $p$-th iterate of $d/dT$ is zero in characteristic $p$, because
its coefficient on every monomial contains a product of $p$
consecutive integers. This also holds for this nonreduced
coefficient ring.

\paragraph{Proof of sharpness.}
In $L(T)$ the exponent of $\pi$ is zero for $r=0,1$ and positive
for every $r\geq2$. Modulo the maximal ideal,
$\overline L(T)=T+T^Q$. In characteristic $p$ its inverse is
\[
 V(T)=\sum_{j\geq0}(-1)^jT^{Q^j},\qquad V(T)+V(T)^Q=T.
\]
Formal inversion commutes with reduction. Hence the
$T^{Q^r}$ coefficient of $U$ is a unit. Since
$E(Z)=\beta U(Z/\beta)$, the coefficient of $Z^{Q^r}$ in $E$
has valuation exactly $-(Q^r-1)/(Q-1)$. For $f=T$ and $a=Q^r$,
$D^aT(0)=a![Z^a]E(Z)$ and $s_p(a)=1$, giving
\[
 v_p(D^aT(0))=
 \frac{a-1}{p-1}-\frac{a-1}{p^2-1}
 =\frac{(a-1)p}{p^2-1}.
\]
This proves all assertions. $\square$

\paragraph{Precise relation to the CM local source.}
The coefficient is
\[
 \frac{p}{p^2-1}=\frac1p+\frac1{p^3}+\frac1{p^5}+\cdots,
\]
exactly the linear slope in Alp\"oge's step~2, p.~3.
For this arithmetic application, now require $\pi$ to be a
uniformizer of the unramified quadratic extension of $\Q_p$.
In that local-field setting the displayed logarithm is a permitted
coordinate for the height-two Lubin--Tate formal group.
The purely formal theorem above allows arbitrary $\pi\in\C_p$
of valuation one and does not assert an integral formal group for
that wider class of parameters.
Bannai--Kobayashi, \emph{Integral structures on $p$-adic Fourier
theory}, Proposition~3.5 and proof, pp.~531--532, first use
$\sum T^{q^r}/\pi^r$ and then treat general integral parameters.
An integral change of parameter with unit linear coefficient
preserves integral coefficient functions. Under a formal-group
isomorphism the normalized invariant derivations differ by a
constant unit, so the valuation statement is unchanged.
In their CM elliptic application, \S5, pp.~543--544, $p$ is odd
inert and the curve has good supersingular reduction; its formal
group is the height-two Lubin--Tate group for
$\pi=-\epsilon p$. The preceding proof allows this unit multiple.

The source provides actual integral functions to which the estimate
applies. On p.~546 it sets, for $b$ prime to $p$,
\[
 \wp_b(z)=\wp(z)-b^2\wp(bz),\qquad
 \phi(t)=\wp_b(\lambda(t))\in\mathcal O_K[[t]],
\]
and constructs an integral modified zeta function as well.
Proposition~5.5, pp.~546--547, identifies the invariant derivatives
and Fourier moments with modified Bernoulli--Hurwitz values.
Theorem~5.1, p.~544, states Katz's stronger floor-divisibility
result for the specified modified values, re-proved in that section.
These are existing CM elliptic applications of the same slope; no
new proof of Katz's full theorem is claimed here.

The slope also equals the valuation of a Lubin--Tate period.
Schneider--Teitelbaum, \emph{$p$-adic Fourier theory},
Lemma~3.4(b), gives
\[
 v_p(\Omega_{\mathrm{LT}})=\frac1{p-1}-\frac1{e(q-1)}.
\]
Putting $e=1$, $q=p^2$ gives $p/(p^2-1)$; the source's appendix
gives the period calculation. The $p$-adic period
$\Omega_{\mathrm{LT}}$ is distinct from the complex period
$\Omega_{\C}$ in the analytic normalization of Alp\"oge's theta
values.

\paragraph{Divided powers give a stronger, testable local identity.}
Katz, \emph{Divisibilities, congruences, and Cartier duality},
Identity~2.4 and Congruences~2.5, p.~670, work in the integral
algebra of invariant differential operators of the coordinate model
$[p](T)=pT+\alpha T^{Q}$, with $\alpha$ a unit and $Q=p^2$.
Define $D(n)$ by
\[
 f(X\mathbin\oplus Y)=\sum_{n\geq0}(D(n)f)(X)Y^n,
 \qquad D=D(1).
\]
Then $D(n)=D^n/n!$ for $n<Q$, and
\[
 D(n)^p\equiv
 \begin{cases}
 0,&p\nmid n,\\
 \alpha^{n/p}D(n/p),&p\mid n
 \end{cases}
 \pmod{p\operatorname{Diff}(G)}.
\]
In particular $D^p\equiv0$ and $D(p)^p\equiv\alpha D$ modulo $p$.
For completeness, the operator-valued translation series
$\Phi(U)=\sum D(n)U^n$ satisfies
$\Phi([p](U))=\Phi(U)^p$ by iterated translation.
Reduction modulo $p$ and comparison of the $U^{pn}$ coefficients
give the congruence. Over characteristic zero,
$\Phi(U)=\exp(D\lambda(U))$; since
$\lambda(U)\equiv U\pmod{U^Q}$, comparison below degree $Q$
gives the divided-power identity. This reconstructs the source's
elementary operator argument and exhibits more integral information
than nilpotence alone.

Katz's Main Theorem~3.11, p.~672, gives the stronger bound
\[
 v_p(D^nf(0))\geq\left\lfloor\frac{np}{p^2-1}\right\rfloor
\]
for integral $f$ satisfying the explicit torsion-trace equation
\[
 \sum_{u\in G[p]}f(T\mathbin\oplus u)=0.
\]
This theorem is cited as background, not reproved here.
In \S~IV, equations~(4.8)--(4.10), p.~676, Katz constructs such an
integral function on the supersingular CM elliptic formal group:
$f=\wp_b-[p]^*\wp_b$. Equations~(4.13)--(4.16), p.~677, identify
its derivatives with the modified Bernoulli--Hurwitz values.
Thus the stronger divisibility has a concrete functional equation
and actual integral elliptic examples. Matching a proposed theta
section to this function, or proving its own torsion-trace identity
in a compatible integral frame, remains a specific additional
comparison.

\paragraph{An integral trace-zero operator and the optimal uniform error one.}
In Katz's exact model $[p](T)=pT+\alpha T^Q$, over the ring $R$
of integers of a finite extension of $\Q_p$, the cited theorem
implies the stronger uniform estimate
\[
 v_p(D^af(0))\geq
       \left\lfloor\frac{ap}{p^2-1}\right\rfloor-1
       \qquad(f\in R[[T]],\ a\geq1).
\]
The constant one is optimal uniformly in $f$ and $a$.
Here is a complete reduction to Katz's theorem. Put
$A=R[[S]]$, $B=R[[T]]$, $S=[p](T)$. Weierstrass division for
\[
 T^Q+(p/\alpha)T-S/\alpha=0
\]
gives a free $A$-basis $1,T,\ldots,T^{Q-1}$ of $B$.
Newton identities yield
\[
 \Tr(1)=Q,\qquad \Tr(T^j)=0\ (1\leq j<Q-1),\qquad
 \Tr(T^{Q-1})=-(Q-1)p/\alpha.
\]
Thus $\Tr_{B/A}(B)=pA$. Define the integral functions
\[
 h=\Tr_{B/A}(f)/p,\qquad
 g=pf-[p]^*h.
\]
Their trace satisfies $\Tr(g)=p\Tr(f)-Qh=0$.
Over the split generic kernel this trace is the sum over translations
by $G[p]$, so $g$ satisfies exactly Katz's torsion-trace equation.
The equality holds integrally as well, since the formal rings are
$p$-torsion free. Katz's theorem therefore applies to $g$.
The invariant derivative satisfies $D[p]^*=p[p]^*D$, and $D$
preserves the integral formal functions. Hence
\[
 pD^af(0)=D^ag(0)+p^aD^ah(0)
\]
has valuation at least $\lfloor ap/(p^2-1)\rfloor$; the second
summand has valuation at least $a$, which is no smaller.
Dividing by $p$ proves the estimate. For $a=Q-1$ and
$f=T^{Q-1}$, one has $D^af(0)=a!$ and
\[
 v_p(a!)=p-1,
 \qquad\left\lfloor\frac{ap}{p^2-1}\right\rfloor=p.
\]
A smaller uniform loss is therefore impossible. $\square$

The argument initially has the exact DVR and coordinate hypotheses
of Katz. It includes the CM elliptic groups described in his \S~IV.
After integral base change from this model, it also holds for
$\mathcal O_{\C_p}$ coefficients: a fixed derivative uses only a
finite jet; approximate its integral coefficients by algebraic
integral coefficients in a finite extension and pass to the closed
valuation ball. Integral changes of formal parameter with unit
linear coefficient preserve the valuation of the normalized
invariant derivative. This is a deduction from Katz's theorem,
without a priority claim. It is not an automatic assertion about
an arbitrary analytic theta normalization.

\paragraph{The remaining theta-frame gate.}
A theta function is a section of a line bundle. Extension as an
invertible sheaf over the formal integral neighborhood gives some
integral frame, but differentiating its coefficient depends on the
frame or connection. The comparison still required is an equality
between the normalized analytic derivatives
$D^a\theta(0)/\Omega_{\C}^a$ and $D_G^af(0)$ for an integral
coefficient function in a specified compatible frame. Existence of
an arbitrary integral trivialization does not prove that equality;
a horizontal or cubical analytic frame need not be integral.

This is an actual obstruction to inferring the derivatives from
the divisor. On the formal characteristic-zero germ
$\C_p[[T]]$, take the trivial line bundle, section one, and frame
$e$. Replace the frame by
$e'=(1+p^{-M}T^a)e$. The two frames agree at the origin and the
section has the same empty divisor, but its coefficient becomes
$f'=(1+p^{-M}T^a)^{-1}$ and
\[
 D^af'(0)=-a!p^{-M}.
\]
Its valuation is arbitrarily negative as $M$ grows. The new frame
is a generic formal frame, not an integral one; no assertion about
invertibility on the whole rigid open unit disc is involved.
Even the exact integral Cartier divisor does not determine higher
jets: for $2\leq a<Q$, all
$f_N(T)=T(1+p^NT^{a-1})$ define the divisor $(T)$ and have first
derivative one, whereas
\[
 D^af_N(0)=a!p^N,
\]
because $E(Z)=Z+O(Z^Q)$. The uniform lower bound remains valid,
but the precise derivatives cannot be recovered from that divisor.

\paragraph{A concrete reduced-theta failure of integrality.}
The distinction occurs for a usual normalized reduced theta.
Take $E:y^2=x^3-x$, $\omega=dx/(2y)$, and $t=-x/y$.
At $p=7$ this is a good supersingular CM elliptic curve.
After replacing $y$ by $2y$, it is precisely the example in
Bannai--Kobayashi, arXiv:math/0610163v4, \S4, pp.~49--50.
Their Example~1.9, pp.~10--11, defines
$\theta(z)=\exp(-e_2^*z^2/2)\sigma(z)$ with $\theta'(0)=1$.
For the square lattice $e_2^*=0$, by its $i$-symmetry, also recorded
in Figure~1, p.~49. Thus $\theta=\sigma$ in this normalization.

The equation $t^2=x/(x^2-1)$ gives
\[
 x=\frac{1+\sqrt{1+4t^4}}{2t^2},\qquad y=-x/t,
 \qquad \omega=\frac{dt}{\sqrt{1+4t^4}}.
\]
Consequently
\[
 \lambda(t)=t-\frac25t^5+\frac23t^9+O(t^{13}).
\]
The Weierstrass equation $\wp''=6\wp^2-2$ gives
$\wp(z)=z^{-2}+z^2/5+z^6/75+O(z^{10})$.
Using $(\log\sigma)''=-\wp$ and the sigma normalization yields
\[
 \log(\sigma(z)/z)=-\frac{z^4}{60}-\frac{z^8}{4200}+O(z^{12}),
 \qquad
 \sigma(z)=z-\frac{z^5}{60}-\frac{z^9}{10080}+O(z^{13}).
\]
Exact composition therefore gives
\[
 \theta(\lambda(t))
 =t-\frac5{12}t^5+\frac{1411}{2016}t^9+O(t^{13}).
\]
Since $2016=7\cdot288$ and $7\nmid1411$, the $t^9$ coefficient
has valuation $-1$ at $7$. This actual normalized reduced theta is
not an integral coefficient function in the stated formal
parameter. The finite calculation makes no assertion about an
asymptotic radius of convergence and does not refute Alp\"oge's
accepted asymptotic theorem. It proves that analytic reducedness
alone does not provide the integral-frame hypothesis needed above.
A separate rational symbolic calculation reproduced these
coefficients; it is not a formal proof certificate.

\paragraph{Framework anchor.}
This finding concerns an integral function, invariant derivation, and
specified local frame.  It reaches the correspondence/polarization
problem of cote~44, scans~7--8, only after a section and integral-model
comparison such as C17 is supplied.  The explicit reduced-theta example
shows why that normalization datum cannot be suppressed.

\subsection{THETA3-C17: the cycle-moduli comparison and its specialization criterion}
\label{C17}

The specialization statement below is a standard consequence of the cited
EGA/SGA functorialities.  The contribution here is the explicit comparison
criterion and identification of the geometric data that an application
to the proposed bridge would have to supply.

\paragraph{The two cycles have different definitions and targets.}
The 032 construction starts with maps from a ball surface to fixed
CM abelian varieties, induced through the Albanese. For
$f:S\to A=\prod_iB_i$, it produces the image-surface cycle
$f_*[S]\in\CH_2(A)$. It is a cycle on the fixed abelian variety
$A$.

By contrast, Li--Zhang \S2.3, p.~11, defines a local special cycle
$\mathcal Z(L)$ by an extension condition on special
quasi-homomorphisms: after applying the framing quasi-isogenies,
each $x\in L$ must extend from the special fibre to an actual
$\mathcal O_k$-linear homomorphism of the indicated
$p$-divisible groups. Globally, \S\S13.2--13.3, p.~82, starts
with the unitary moduli object $(A_0,A)$, its actions,
polarizations, and level, and considers tuples of special
homomorphisms in the allowed prime-to-$p$ isogeny category with
prescribed Hermitian Gram matrix $T$ and prescribed finite test
function. This is a moduli condition on a unitary Shimura model.
It is not supplied by an arbitrary map from its surface to fixed
abelian varieties. In particular the two definitions neither have
the same ambient variety nor automatically the same codimension.

A literal comparison therefore requires a specified algebraic
correspondence from the relevant Shimura model to the fixed target
$A$, a theorem identifying its action on the selected special cycle
with the projected image-surface cycle, and integral extensions
compatible with the desired local model. This is an identification
gate, not an assertion that such a correspondence is impossible.
C04 still supplies the separate rank requirement on the target
Hodge summand.

\paragraph{Exact primary EGA and SGA locators.}
EGA~IV$_3$, Theorem~8.8.2, pp.~28--29, and Scholie~8.8.3,
pp.~33--34, descend finite-presentation schemes, morphisms, and
finite collections of relations through filtered limits.
EGA~IV$_2$, Theorem~6.9.1, p.~153, is generic flatness. These give
spreading of finite specified geometric data and flatness after
shrinking a suitable integral base. They do not produce a desired
cycle, and shrinking does not ensure good behavior at an arbitrarily
specified bad prime.

For specialization use the original SGA~4, tome~III,
Expos\'e~XVI, Theorem~2.1, p.~213, and Corollary~2.3, p.~215,
with locally constant prime-to-residue-characteristic coefficients.
The locally constant hypothesis is explicit in those statements;
one need not rely on the misprinted formulation of Corollary~2.2.
For cycle classes use SGA~4$\frac12$, \emph{La classe de cohomologie
associ\'ee \`a un cycle}: Definition~2.3.2, p.~145;
the perfect-weight construction~2.3.3, p.~146;
Lemma~2.3.7 and Theorem~2.3.8, p.~148; and
Remark~2.3.9, p.~150. Lemma~2.3.7 gives base-change compatibility.
Theorem~2.3.8 gives pullback and proper-intersection compatibility
under its fibrewise codimension and perfect-weight hypotheses.
Remark~2.3.9 explicitly uses alternating sums of Tor for proper
intersection multiplicities. These original scans were inspected;
none of these source theorems is claimed to have been reproved here.

\paragraph{A precise cohomological criterion.}
Let $R$ be a strictly henselian discrete valuation ring,
$\ell$ invertible in $R$, and $X,Y$ smooth proper over $R$.
Let $\Gamma\subset X\times_RY$ be a relative flat cycle of pure
codimension $c$ with the preceding cycle-class hypotheses; for
example its structure sheaf is a perfect weight on the regular
ambient product. Fix geometric generic and special fibres and the
specialization isomorphisms. Then the action of $\Gamma$ commutes
with specialization:
\[
 \operatorname{sp}_Y\circ(\Gamma_{\bar\eta})_*
       =(\Gamma_{\bar s})_*\circ\operatorname{sp}_X.
\]
The correspondence action has the usual degree shift
$2c-2\dim(X/R)$ and Tate twist $c-\dim(X/R)$.

\paragraph{Proof from the cited functorialities.}
Write the action as pullback to the product, cup product with
$\operatorname{cl}(\Gamma)$, and the proper smooth projection's
trace. The cited specialization theorem identifies the fibre
cohomologies. Pullback and cup product commute with base change;
Lemma~2.3.7 identifies the two restrictions of the cycle class;
and the trace functoriality in the construction~2.3.1--2.3.4
commutes with that base change. Combining the three operations
proves the identity. Flatness makes the fibre restriction of
$\mathcal O_\Gamma$ agree with its derived base change, so it is
the stipulated fibre cycle class. $\square$

Consequently, once actual correspondences and their relative cycle
classes are provided, their cohomological action and any certified
rank calculation can be transported through this comparison. The
criterion neither lifts arbitrary special-fibre cycles nor proves
that an arbitrary Hodge projector is algebraic. The same
cycle-class/intersection compatibility shows that homological
triviality implies numerical triviality: pair a zero cycle class
with each complementary algebraic class and apply the trace.
The converse is not supplied by EGA/SGA or by the local
intersection formula. Passing among a geometric correspondence,
its homological realization, and a desired Hodge tensor therefore
requires the explicit maps and rank statements described above.

\paragraph{Archival anchor.}
Cote~16, scan~43, states the complementary-dimension issue for algebraic
classes; scans~44--45 formulate cohomological inverse maps induced by
algebraic correspondences; scan~46 continues the discussion of
specialization and hyperplane reductions.  C17 concerns the functoriality
of already specified relative correspondences.  It does not establish the
existence of the conjectural correspondences in that letter.  Cote~44,
scans~7--8, supplies the related Albanese and polarization interpretation.
